\section{Additional ablation and discussion}
\label{sec:ablation_discussion}
\subsection{Comparison of compact 3D reconstruction strategies}
\begin{table}[t]
\centering
\caption{\textbf{Comparison of novel view synthesis with compact Gaussians generation strategies.} }
\label{suptab:coarse_reconstruction}
\resizebox{0.75\linewidth}{!}{
\begin{tabular}{l|cccc}
\toprule
\multirow{2}{*}{Methods}  & \multicolumn{4}{c}{Average} \\
&  PSNR$\uparrow$ & SSIM$\uparrow$ & LPIPS$\downarrow$ & \#G$\downarrow$ \\
\midrule
Sampling & 21.340& 0.665& 0.272& \textbf{2K} \\
Voxelize & 19.957& 0.609& 0.403& 4K \\
Ours & \textbf{22.387}& \textbf{0.713}& \textbf{0.259}& \textbf{2K} \\
\bottomrule
\end{tabular}}
\end{table}
To validate the effectiveness of our compact reconstruction strategy, we conduct ablation studies comparing different approaches to obtain compact representations. For the sampling baseline, we first estimate per-pixel Gaussian centers and their corresponding Gaussians, then downsample by 8$\times$ to predict only 2,048 Gaussians, matching the number of Gaussians with our method. For the voxelization baseline, we follow AnySplat's~\cite{jiang2025anysplat} strategy: we first estimate per-pixel Gaussians, then voxelize them with additional layers. Since AnySplat's original voxel size cannot sufficiently reduce the number of Gaussians, we increase the voxel size to 0.2, which results in approximately 4K Gaussians to best match the 2,048 Gaussians produced by our method.

As shown in Tab.~\ref{suptab:coarse_reconstruction}, our strategy demonstrates superior performance compared to sampling or voxelization methods while using fewer Gaussians than voxelization. Both sampling and voxelization rely on per-pixel estimation, which restricts them to pixel locations and prevents accurate estimation of necessary regions. Additionally, voxelization requires heuristic determination of voxel size and can reduce details in high-frequency regions. In contrast, our approach does not depend on pixel locations as learnable tokens can flexibly estimate Gaussians at appropriate locations, enabling more effective compact Gaussian field reconstruction.

\subsection{Additional ablation for feature decoder}
We additionally ablate the components of \oursF. In Tab.~\ref{suptab:addi_abl_comp3df}, we conduct experiments on (1) detaching the copied attention weights from \ours, and (2) propagating feature loss to Gaussian geometry attributes (means, covariances, spherical harmonics, and opacities). When feature loss propagates to the copied attention weights, \ours\ receives ambiguous gradients because the features are not perfectly multi-view consistent. Consequently, the attention mechanism cannot correctly identify correspondences between learnable query tokens and the encoder features $\mathcal{E}(\cdot)$~(e.g., VGGT) for \ours.
To propagate feature loss to Gaussian attributes, we modify the CUDA rasterizer from Feature-3DGS~\cite{zhou2024feature}, which originally propagates feature loss only to Gaussian feature attributes. We extend the CUDA rasterizer to propagate feature loss to all Gaussian attributes. However, this also degrades geometry estimation results because foundation model features are not perfectly multi-view consistent.
We hypothesize that with perfectly multi-view invariant features, feature loss could improve reconstruction quality, especially since photometric loss is also imperfect, as real-world RGB images contain visual artifacts such as appearance variations, lighting changes, and noise.

\begin{table}[t]
\centering
\caption{\textbf{Additional ablation studies for \oursF.} }
\label{suptab:addi_abl_comp3df}
\resizebox{\linewidth}{!}{
\begin{tabular}{cc|ccccc}
\toprule
\multirow{2}{*}{\shortstack{Detach\\attn.}} & \multirow{2}{*}{\shortstack{Detach\\ $\mathcal{L}_\text{feat}$ to $G_i$} } & \multicolumn{5}{c}{Lseg~\cite{li2022language}} \\
& & mIOU$\uparrow$ & Acc.$\uparrow$ & PSNR$\uparrow$ & SSIM$\uparrow$ & LPIPS$\downarrow$ \\
\midrule
\ding{55} & \ding{51} & 0.490& 0.761& 23.078& 0.750&0.293\\
\ding{51} & \ding{55} & 0.490& 0.757& 23.243& 0.754&0.286\\
\textbf{\ding{51}} & \textbf{\ding{51}} & \textbf{0.513}& \textbf{0.783}& \textbf{23.886}& \textbf{0.770}& \textbf{0.285}\\
\bottomrule
\end{tabular}}
\end{table}

\subsection{Discussion of autoencoder in existing methods}
\begin{figure}[t]
    \centering
    \includegraphics[width=\linewidth, height=0.7\linewidth]{supple_sec/supple_figure/assets/sup_loss.pdf}
    \vspace{-15pt}
    \caption{\textbf{Convergence speed improvement by eliminating autoencoder in our framework.}} 
    \label{supfig:convergence}\vspace{-10pt}
\end{figure}

To reduce the memory consumption of Gaussians, prior works~\cite{fan2024large, lee2025cf3} commonly adopt autoencoders to compress Gaussian feature dimensions and restore them via a decoder. However, autoencoder-based compression inevitably introduces information loss.
Moreover, as shown in Fig.~\ref{supfig:convergence}, removing the autoencoder allows our model to converge \textbf{7.1$\times$} faster. With an autoencoder framework, additional training time is required for the encoder–decoder architecture. In contrast, without an autoencoder, we directly leverage the attention mechanism of \ours\ to aggregate and store features. Since the attention maps in \ours\ already identify the necessary features for each learnable token, each feature token naturally attends to the same multi-view regions as its corresponding Gaussian token, effectively reusing learned correspondences for feature aggregation. This emergent property enables highly efficient feature lifting with minimal additional training overhead.

\subsection{Discussion of FPS}
As shown in Tab.~\textcolor[HTML]{367DBD}{3} and Tab.~\textcolor[HTML]{367DBD}{4}, our method reduces the number of Gaussians by \textbf{65$\times$} (only \textbf{2K}) compared to per-pixel approaches~\cite{fan2024large,ye2024no}, yet the FPS does not increase linearly. This is due to GPU saturation: modern accelerators (e.g., H100) are highly parallelized, and the rasterization overhead for small primitive counts (e.g., 2K) is dominated by fixed pipeline latencies rather than the sorting or splatting cost. Consequently, both LSM and our method operate within the saturation regime of the hardware. However, while the FPS gain saturates, our \textbf{15$\times$} reduction in memory (4.1MB vs. 61.5MB) provides a critical advantage for bandwidth-constrained applications where storage is the primary bottleneck, not rendering speed.

\subsection{Limitations}
Although our view-invariant feature decoder, \oursF, allows lifting arbitrary 2D features into 3D, we have not evaluated all recent foundation models. For instance, we believe that leveraging features from the Segment Anything Model~(SAM)~\cite{kirillov2023segment} could enable robust multi-view consistent segmentation if the features are aggregated within our framework before being decoded by a pre-trained SAM decoder; however, we did not analyze these specific features in this work.

In addition, following prior works, our experimental validation primarily focuses on multi-view open-vocabulary segmentation. To achieve holistic 3D scene understanding, future work could explore integrating our feature fields with recent Multimodal Large Language Models~(MLLMs)~\cite{liu2023visual,liu2024improved} to enable 3D Scene Question Answering. Furthermore, since our framework enables novel view feature rendering without information loss or compression artifacts, it holds significant promise for integration with Vision-Language-Action~(VLA) models or robotics applications. However, such scenarios are frequently dynamic, whereas our current framework is limited to static scene reconstruction. Extending our compact representation to dynamic scenes remains an exciting avenue that would unlock potential across various autonomous fields.
%\clearpage
